\newif\ifglobal

%%%%%%%%%%%%%%%%%%%%%%%
%%% Compile to view %%%
%%%%%%%%%%%%%%%%%%%%%%%

%%%%%%%%%%%%%%%%%%%%%%%%%%%%%%%%%%%%%%%%%%%%%%%%%%%
%%% Readme:                                     %%%
%%% https://www.overleaf.com/read/bwdjvfjksvjy/ %%%
%%%%%%%%%%%%%%%%%%%%%%%%%%%%%%%%%%%%%%%%%%%%%%%%%%%

% >>> M?: marks a module setting number or important notice

% >>> M4: Set {./relative-path} to external files for this module <<<
\def \modulefiles {./22-RNG}
% To add an external file, use:
% (Replace '---' with actual filename)
% '\input{\modulefiles/tables/---.tex}' for tables
% '\includegraphics{\modulefiles/figures/---}' for figures
% '\subfile{\modulefiles/---}' for register files

% >>> M1: This module is in global project? <<<
% 'YES' - uncomment; 'NO' - comment out
\globaltrue

\ifglobal
    \documentclass[main\_\_EN.tex]{subfiles}

\else
    \documentclass[openany, 10pt]{book}

    % >>> M2: In ./chip-spec-settings/preamble-trm-module, also find
    %         settings for visibility of todo notes, watermark <<<
    \usepackage{preamble-trm-module}

    % >>> M3: Temporary labels in English,
    %         you can add your own labels there <<<
    \myexternaldocument{temp-labels-en}

\fi %\ifglobal


\setcounter{secnumdepth}{4}




% >>> M5: If this module needs any updates in
% './chip-spec-settings/preamble-trm-module.sty'
% (include additional package, etc.), contact your module PM
% or create the file './chip-spec-settings/changelog.tex' make the
% necessary updates yourself and document them in this file <<<

\raggedright


\begin{document}


% Sets language into which pre-defined bits of text are translated
\selectlanguage{english}

% Do not delete this block
\ifglobal
% Add the following front matter if
% this module is outside of global project
\else
    \InputIfFileExists{readme}{}{}
    {\let\clearpage\relax \listoftodos}
    {\let\clearpage\relax \listofdonetodos}
    \tableofcontents
    %\listoftables
    %\listoffigures
\fi


%%%%%%%%%%%%%%%%%%%%%%%%%
%%% TRM Module Begins %%%
%%%%%%%%%%%%%%%%%%%%%%%%%

\hypertarget{rng}{}
\chapter{Random Number Generator (RNG)}
\label{mod:rng}


\section{Introduction}

The ESP32-S2 contains a true random number generator, which generates 32-bit random numbers that can be used for cryptographical operations, among other things.

\section{Features}

The random number generator generates true random numbers, which means random number generated from a physical process, rather than by means of an algorithm. No number generated within the specified range is more or less likely to appear than any other number.

\section{Functional Description}\label{sec:rng-func-descr}

Every 32-bit value that the system reads from the \hyperref[regdesc:RNGDATAREG]{RNG\_DATA\_REG} register of the random number generator is a true random number. These true random numbers are generated based on the thermal noise in the system and the asynchronous clock mismatch.

Thermal noise comes from the high-speed ADC or SAR ADC or both. Whenever the high-speed ADC or SAR ADC is enabled, bit streams will be generated and fed into the random number generator through an XOR logic gate as random seeds.

When the RC\_FAST\_CLK clock is enabled for the digital core, the random number generator will also sample RC\_FAST\_CLK (8 MHz) as a random bit seed. RC\_FAST\_CLK is an asynchronous clock source and it increases the RNG entropy by introducing circuit metastability. However, to ensure maximum entropy, it’s recommended to always enable an ADC source as well.

\begin{figure}[H]
    \centering
    \includegraphics[width=0.8\textwidth]{\modulefiles/figures/06-RNG-noice-source}
    \caption{Noise Source}
    \label{figure:RNG_Noice_Source}
\end{figure}

When there is noise coming from the SAR ADC, the random number generator is fed with a 2-bit entropy in one clock cycle of RC\_FAST\_CLK (8 MHz), which is generated from an internal RC oscillator (see Chapter \ref{mod:resclk} \textit{\nameref{mod:resclk}} for details). Thus, it is advisable to read the \hyperref[regdesc:RNGDATAREG]{RNG\_DATA\_REG} register at a maximum rate of 500 kHz to obtain the maximum entropy.

When there is noise coming from the high-speed ADC, the random number generator is fed with a 2-bit entropy in one APB clock cycle, which is normally 80 MHz. Thus, it is advisable to read the \hyperref[regdesc:RNGDATAREG]{RNG\_DATA\_REG}
register at a maximum rate of 5 MHz to obtain the maximum entropy.



\section{Programming Procedure}\label{rng-programming}

When using the random number generator, make sure at least either the SAR ADC, high-speed ADC, or RC\_FAST\_CLK  is enabled. Otherwise, pseudo-random numbers will be returned.
\begin{itemize}
    \item SAR ADC can be enabled by using the DIG ADC controller. For details, please refer to Chapter \ref{mod:sensor} \textit{\nameref{mod:sensor}}.
    \item High-speed ADC is enabled automatically when the Wi-Fi module is enabled.
    \item RC\_FAST\_CLK is enabled by setting the \hyperref[fielddesc:RTCCNTLDIGCLK8MEN]{RTC\_CNTL\_DIG\_CLK8M\_EN} bit in the \hyperref[regdesc:RTCCNTLCLKCONFREG]{RTC\_CNTL\_CLK\_CONF\_REG} register.
\end{itemize}

\begin{tiplisting}
Note that, when the Wi-Fi module is enabled, the value read from the high-speed ADC can be saturated in some extreme cases, which lowers the entropy. Thus, it is advisable to also enable the SAR ADC as the noise source for the random number generator for such cases.
\end{tiplisting}

When using the random number generator, read the \hyperref[regdesc:RNGDATAREG]{RNG\_DATA\_REG} register multiple times until sufficient random numbers have been generated. Ensure the rate at which the register is read does not exceed the frequencies described in section \ref{sec:rng-func-descr} above.

\section{Base Address } \label{Label:RNG_Base_Address}

Users can access the random number generator with two base addresses, which can be seen in Table \ref{table:RNG_ENTRY}. For more information about accessing peripherals from different buses, please see Chapter \ref{mod:sysmem} \textit{\nameref{mod:sysmem}}.\\

\begin{longtable}[c]{ | l | C{3cm} | }
\caption{Random Number Generator Base Address} \label{table:RNG_ENTRY}
\\\hline
\rowcolor{lightgray}
    Bus to Access Peripheral  & Base Address\\
    \hline
    PeriBUS1      & \hexprefix{}3F435000      \\
    \hline
    PeriBUS2      & \hexprefix{}60035000      \\
    \hline
\end{longtable}


\newpage
\hypertarget{rng-reg-summ}{}
\section{Register Summary}

The addresses in the following table are relative to the random number generator base addresses provided in Section \ref{Label:RNG_Base_Address}.
\begin{longtable}{ | p{5cm} | p{6cm} | p{3cm} | p{1cm} | }
\hline\rowcolor{lightgray}
Name & Description & Address & Access \\ \hline
\hyperref[regdesc:RNGDATAREG]{RNG\_DATA\_REG} & Random number data & \hexprefix{}0110 & RO \\ \hline
\end{longtable}

\section{Register}
The address in this section is relative to the random number generator base addresses provided in Section \ref{Label:RNG_Base_Address}.

\begin{register}{H}{RNG\_DATA\_REG}{\hexprefix{}0110}\label{regdesc:RNGDATAREG}
\regfield{}{32}{0}{{\hexprefix{}00000000}}%
\reglabel{Reset}\regnewline%
\begin{regdesc}\begin{reglist}
\item [RNG\_DATA] Random number source. (RO)
\end{reglist}\end{regdesc}
\end{register}

\end{document}
